\section{Introduction}
\label{sec:intro}
Evidence for cognitive behavioral therapy (CBT) in adult attention-deficit/hyperactivity disorder (ADHD) warrants a question more precise than whether the treatment ``works.''
A Cochrane review identified possible short-term benefits but emphasized limitations in evidence quality and the need for longer-term evidence \citep{Lopez2018}.
A randomized trial in medication-treated adults with persistent symptoms found greater improvement with CBT than with relaxation and educational support \citep{Safren2010}.
That result concerns a particular population and treatment comparison; the present review asks whether baseline executive function modifies such treatment contrasts.
Recent meta-analyses also examine executive-function outcomes and everyday functioning, providing reasons to distinguish what improves from who benefits \citep{Liu2026,LopezPinar2026}.

The author's archived blog and prior critical review motivate a narrower hypothesis: cognitive restructuring may be harder to use when executive resources are limited.
We treat this as a proposed explanation to test, not as evidence that CBT is ineffective.
Related pilot studies in older adults with anxiety or depression examine executive functioning and cognitive flexibility, but their populations and analyses do not directly answer the adult-ADHD moderation question \citep{Mohlman2005,Johnco2014}.
In particular, \citet{Johnco2014} did not find that pretreatment flexibility predicted overall treatment outcome, although associations with aspects of restructuring skill were reported.
This mixed pattern cautions against turning a plausible cognitive-demand account into a settled clinical conclusion.

The distinction between a predictor, moderator and outcome is central.
A baseline characteristic may predict subsequent improvement without changing the effect of CBT relative to another treatment; moderation concerns that difference in treatment effects \citep{Kraemer2002}.
Similarly, evidence that executive function improves following CBT does not by itself test baseline executive function as a moderator; the outcome and interaction questions require different analyses \citep{Liu2026,Kraemer2002}.
We will therefore extract the timing of executive-function measurement, the treatment contrast and the actual statistical model, rather than treating every mention of executive function and outcome as a moderator result.
We will also distinguish multicomponent CBT from isolated cognitive restructuring, so that a component critique is evaluated at the level for which evidence exists.

This project contributes a prospective protocol and a reproducible search baseline for a critical meta-review, not a completed efficacy verdict.
The primary synthesis will map systematic reviews of adult-ADHD CBT; a supplementary synthesis will identify original studies that can address baseline executive-function prediction or moderation.
Search and selection reporting will follow the Preferred Reporting Items for Systematic Reviews and Meta-Analyses (PRISMA) framework and its search extension \citep{Page2021,Rethlefsen2021}.
The review is designed to retain supportive, contradictory and null evidence and to separate absence of a direct test from evidence of no moderation.
